\section*{Hoofdstuk \ref{ch:g11:fossil-fuel-geology} --- De geologie van steenkool, olie en gas}

\begin{solution}{exo:g11:fossil-fuel-geology:1}
Dood plankton bezinkt op een zuurstofarme zeebodem; het organische materiaal
wordt \omterm{def:g11:fossil-fuel-geology:source-rock}{kerogeen} in een modder; begraving onder jongere sedimenten; \omterm{def:g11:fossil-fuel-geology:oil-window}{rijping} in het
\omterm{def:g11:fossil-fuel-geology:oil-window}{olievenster}; \omterm{def:g11:fossil-fuel-geology:migration}{migratie}; olie verzamelt zich onder een \omterm{def:g11:fossil-fuel-geology:reservoir}{afdekgesteente}.
\end{solution}

\begin{solution}{exo:g11:fossil-fuel-geology:2}
(1) \omterm{def:g11:fossil-fuel-geology:source-rock}{Brongesteente}. (2) \omterm{def:g11:fossil-fuel-geology:reservoir}{Reservoirgesteente}. (3) \omterm{def:g11:fossil-fuel-geology:reservoir}{Afdekgesteente}. (4) \omterm{def:g11:fossil-fuel-geology:reservoir}{Afdekgesteente}. (5) \omterm{def:g11:fossil-fuel-geology:reservoir}{Reservoirgesteente}.
\end{solution}

\begin{solution}{exo:g11:fossil-fuel-geology:3}
Bovenkant: $(60 - 10)/25 = \qty{2.0}{km}$; onderkant: $(120 - 10)/25 = \qty{4.4}{km}$.
Het venster is $60/25 = \qty{2.4}{km}$ dik.
\end{solution}

\begin{solution}{exo:g11:fossil-fuel-geology:4}
(1) Veen, bruinkool, \omterm{def:g11:fossil-fuel-geology:coal}{bitumineuze kool}, antraciet. (2) Antraciet, 86 à
\qty{97}{\percent} koolstof. (3) Antraciet, voorbij ongeveer \qty{5}{km} en
\qty{150}{\degreeCelsius}.
\end{solution}

\begin{solution}{exo:g11:fossil-fuel-geology:5}
(1) De vloeistoffen sorteren zich naar dichtheid: gas is het lichtst, dan olie,
dan water. (2) Olie en gas stijgen omdat ze lichter zijn dan water; alleen
een gesteente waar ze niet doorheen kunnen, houdt ze tegen. Kleisteen (of schalie) en zout.
\end{solution}

\begin{solution}{exo:g11:fossil-fuel-geology:6}
(1) A: $(60 - 15)/20 = \qty{2.25}{km}$ en $(120 - 15)/20 = \qty{5.25}{km}$.
B: $45/40 \approx \qty{1.13}{km}$ en $105/40 \approx \qty{2.63}{km}$.
(2) Op \qty{2}{km}: in A $15 + 40 = \qty{55}{\degreeCelsius}$, nog niet
rijp, geen olie; in B $15 + 80 = \qty{95}{\degreeCelsius}$, in het \omterm{def:g11:fossil-fuel-geology:oil-window}{olievenster}:
het maakt olie. (3) In bekken A: daar ontstaat olie pas onder
\qty{2.25}{km}.
\end{solution}

\begin{solution}{exo:g11:fossil-fuel-geology:7}
(1) $11 \times 3 = \qty{33}{m}$ veen. (2) $\qty{33000}{mm}$ met
\qty{1}{mm} per jaar: ongeveer \num{33000} jaar. (3) Het veen blijft alleen bewaard waar
het onder stilstaand water ligt, buiten bereik van zuurstof; terwijl het zich ophoopt,
moet de bodem in hetzelfde tempo dalen om het oppervlak drassig te houden.
\end{solution}

\begin{solution}{exo:g11:fossil-fuel-geology:8}
(1) $20 + 30 \times 1.2 = \qty{56}{\degreeCelsius}$;
$20 + 30 \times 2.3 = \qty{89}{\degreeCelsius}$;
$20 + 30 \times 4.0 = \qty{140}{\degreeCelsius}$.
(2) 120 miljoen jaar geleden onrijp; 95 miljoen jaar geleden in het
\omterm{def:g11:fossil-fuel-geology:oil-window}{olievenster}; 20 miljoen jaar geleden in het \omterm{def:g11:fossil-fuel-geology:oil-window}{gasvenster}.
\end{solution}

\begin{solution}{exo:g11:fossil-fuel-geology:9}
(1) $2580 - 2400 = \qty{180}{m}$. (2) Gas vult
$2450 - 2400 = \qty{50}{m}$, olie $2580 - 2450 = \qty{130}{m}$. (3) Ze lopen
weg bij het overlooppunt, onder de diepste gesloten hoogtelijn, en migreren verder
omhoog, naar een andere \omterm{def:g11:fossil-fuel-geology:trap}{val} of naar het oppervlak.
\end{solution}

\begin{solution}{exo:g11:fossil-fuel-geology:10}
(1) Bovenkant van het venster op \qty{3}{km}: $(60 - 20)/3 \approx
\qty{13}{\degreeCelsius/km}$; onderkant: $(120 - 20)/3 \approx
\qty{33}{\degreeCelsius/km}$. (2) Op \qty{3}{km} is de temperatuur
$20 + 75 = \qty{95}{\degreeCelsius}$ tot $20 + 102 = \qty{122}{\degreeCelsius}$:
in het \omterm{def:g11:fossil-fuel-geology:oil-window}{olievenster}, of aan de onderkant ervan bij de steilste gradiënten. De regel is
redelijk: in de meeste bekkens is een \omterm{def:g11:fossil-fuel-geology:source-rock}{brongesteente} onder \qty{3}{km} gesteenten
begonnen olie te maken.
\end{solution}

\begin{solution}{exo:g11:fossil-fuel-geology:11}
(1) \qty{30}{\percent}: bruinkool; \qty{70}{\percent}: \omterm{def:g11:fossil-fuel-geology:coal}{bitumineuze kool};
\qty{92}{\percent}: antraciet. (2) Het antraciet, voorbij ongeveer
\qty{5}{km} en \qty{150}{\degreeCelsius}. (3) Zwart en glazig.
\end{solution}

\begin{solution}{exo:g11:fossil-fuel-geology:12}
(1) $1100/25 = \qty{44}{m}$ per miljoen jaar. (2) $0.044 \times 30
\approx \qty{1.3}{\degreeCelsius}$ per miljoen jaar. (3) Het venster beslaat
\qty{60}{\degreeCelsius}: $60/1.32 \approx 45$ miljoen jaar. (4) De begraving
ging na 95 miljoen jaar geleden trager (ongeveer 24, daarna 23 en
\qty{20}{m} per miljoen jaar): het \omterm{def:g11:fossil-fuel-geology:source-rock}{brongesteente} warmde langzamer op.
\end{solution}

\begin{solution}{exo:g11:fossil-fuel-geology:13}
(1) $(80 - 51)/(117 - 51) = 29/66 \approx 0.44$: ongeveer \qty{44}{\percent}.
(2) Ze migreerde omhoog voordat er erboven een \omterm{def:g11:fossil-fuel-geology:trap}{val} bestond: ze ontsnapte naar het
oppervlak, waar ze werd geoxideerd, of verzamelde zich in een andere \omterm{def:g11:fossil-fuel-geology:trap}{val}. (3) Rond 40
miljoen jaar geleden had het \omterm{def:g11:fossil-fuel-geology:source-rock}{brongesteente} het \omterm{def:g11:fossil-fuel-geology:oil-window}{olievenster} verlaten: het maakte alleen gas,
en de anticlinaal kon er alleen gas van bevatten.
\end{solution}

\begin{solution}{exo:g11:fossil-fuel-geology:14}
(1) Bovenkant: $(60 - 20)/40 = \qty{1.0}{km}$; onderkant: $(120 - 20)/40 =
\qty{2.5}{km}$. (2) \qty{1.0}{km} ligt tussen \qty{0.4}{km} (150 miljoen
jaar geleden) en \qty{1.2}{km} (120): $150 - 30 \times 0.6/0.8 = 127.5$
miljoen jaar geleden. \qty{2.5}{km} ligt tussen \qty{2.3}{km} (95) en
\qty{3.0}{km} (66): $95 - 29 \times 0.2/0.7 \approx 86.7$ miljoen jaar
geleden. (3) Geen olie: het \omterm{def:g11:fossil-fuel-geology:source-rock}{brongesteente} verliet het \omterm{def:g11:fossil-fuel-geology:oil-window}{olievenster} voordat de anticlinaal
ontstond, en zijn olie was weggemigreerd. Hij zou gas kunnen bevatten dat na 80
miljoen jaar geleden ontstond.
\end{solution}

\begin{solution}{exo:g11:fossil-fuel-geology:15}
(1) Minstens \qty{5}{km}. (2) $15 + 30 \times 5 = \qty{165}{\degreeCelsius}$,
boven \qty{150}{\degreeCelsius}: dat klopt. (3) Minstens \qty{5}{km}
gesteente is opgeheven en weggeërodeerd. $\qty{5e6}{mm}$ in hoogstens
\num{299} miljoen jaar geeft minstens $5 \times 10^6/2.99 \times 10^8
\approx \qty{0.017}{mm}$ per jaar gemiddeld.
\end{solution}

\begin{solution}{pb:g11:fossil-fuel-geology:1}
\textbf{1.} \omterm{def:g11:fossil-fuel-geology:source-rock}{Brongesteente}: de zwarte schalie; reservoir: de poreuze zandsteen;
\omterm{def:g11:fossil-fuel-geology:reservoir}{afdekgesteente}: de kleisteen erboven.

\textbf{2.} Plankton dat op een zuurstofarme zeebodem bezonk, ontsnapte aan
afbraak en werd in de modder begraven.

\textbf{3.} Het \omterm{def:g11:fossil-fuel-geology:source-rock}{brongesteente} in het Jura (165 ligt tussen 201.4 en
143.1). De zandsteen begon in het Jura en werd grotendeels afgezet in het
Vroeg-Krijt.

\textbf{4.} Het organische materiaal dat verspreid zit in een sedimentair gesteente, onoplosbaar in
organische oplosmiddelen. Het maakt pas olie wanneer het tot ongeveer
\qty{60}{\degreeCelsius} is verhit, en daarvoor is begraving nodig.

\textbf{5.} $400 + 800 + 1100 + 700 + 700 + 300 - 500 = \qty{3500}{m}$.

\textbf{6.} \qty{0.4}{km}, \qty{1.2}{km}, \qty{2.3}{km}, \qty{3.0}{km},
\qty{3.7}{km}, \qty{4.0}{km}; vandaag \qty{3.5}{km}.

\textbf{7.} $400/15 \approx 26.7$; $800/30 \approx 26.7$; $1100/25 = 44$;
$700/29 \approx 24.1$; $700/31 \approx 22.6$; $300/15 = 20$ (\unit{m} per
miljoen jaar).

\textbf{8.} 20 miljoen jaar geleden, op \qty{4.0}{km}.

\textbf{9.} Het bekken werd opgeheven en \qty{500}{m} gesteente werd van de
bovenkant weggeërodeerd.

\textbf{10.} De basis van de zandsteen werd 150 miljoen jaar geleden afgezet, boven
op \qty{400}{m} schalie en kleisteen: $3.5 - 0.4 = \qty{3.1}{km}$.

\textbf{11.} $T = 20 + 30 z$: \qty{32}{\degreeCelsius}, \num{56}, \num{89},
\num{110}, \num{131}, \qty{140}{\degreeCelsius}; vandaag
\qty{125}{\degreeCelsius}.

\textbf{12.} $40/30 \approx \qty{1.33}{km}$ en $100/30 \approx
\qty{3.33}{km}$.

\textbf{13.} Tussen 120 miljoen jaar geleden (\qty{56}{\degreeCelsius}) en
95 miljoen jaar geleden (\qty{89}{\degreeCelsius}).

\textbf{14.} De temperatuur steeg 33 graden in 25 miljoen jaar; het duurde
$25 \times 4/33 \approx 3.0$ miljoen jaar om de 4 graden van
\num{56} naar \qty{60}{\degreeCelsius} te stijgen: ongeveer $120 - 3.0 = 117$ miljoen jaar
geleden.

\textbf{15.} Tussen 66 miljoen jaar geleden (\qty{110}{\degreeCelsius}) en
35 (\qty{131}{\degreeCelsius}): $66 - 31 \times 10/21 \approx 51.2$ miljoen
jaar geleden.

\textbf{16.} \qty{140}{\degreeCelsius}, 20 miljoen jaar geleden. Met
\qty{125}{\degreeCelsius} ligt het vandaag in het \omterm{def:g11:fossil-fuel-geology:oil-window}{gasvenster}: het maakt vooral
methaan.

\textbf{17.} Zonder \omterm{def:g11:fossil-fuel-geology:trap}{val} erboven migreerde ze omhoog en ontsnapte ze, naar het
oppervlak of naar andere vallen.

\textbf{18.} $(80 - 51.2)/(117.0 - 51.2) \approx 0.44$.

\textbf{19.} Gas bovenaan, dan olie, dan water.

\textbf{20.} Het \omterm{def:g11:fossil-fuel-geology:source-rock}{brongesteente} kwam ongeveer \textbf{117
miljoen jaar geleden} in het \omterm{def:g11:fossil-fuel-geology:oil-window}{olievenster}.
\end{solution}
